\section{Problem Statement and Architectural Pipeline}

\subsection{Operational Problem Context}
In real-world applications of Comparative Study of Classification Algorithms for Early Disease Risk Prediction, systems encounter Severe class imbalance (cardiac event rate $< 9\%$), noisy physiological measurements, high dimensional interaction effects between metabolic indicators, and varying computational complexity.. The operational workflow requires processing raw measurements, mitigating data artifacts, orchestrating model selection, and validating predictions under rigorous cross-validation paradigms.

\subsection{System Pipeline Architecture}
Figure~\ref{fig:pipeline} illustrates the end-to-end computational pipeline developed to solve the target machine learning challenge. The architecture progresses from ingestion and preprocessing through algorithmic estimation, hyperparameter tuning, and metric verification.

\begin{figure}[H]
\centering
\resizebox{\textwidth}{!}{\begin{tikzpicture}[
    block/.style={rectangle, draw=black!80, fill=blue!8, thick, rounded corners, minimum height=1.1cm, text width=2.5cm, font=\footnotesize},
    arr/.style={-{Stealth[length=2mm]}, thick}
]
\node[block, align=center] (n1) at (0,0) {1. EHR Cohort ($N=14k$)\\Cardiovascular Lab Data};
\node[block, align=center] (n2) at (3.2,0) {2. Preprocessing Pipeline\\KNN Impute \& SMOTE};
\node[block, align=center] (n3) at (6.4,0) {3. Feature Scaling\\Robust Z-Score Normalization};
\node[block, align=center] (n4) at (9.6,0) {4. Model A: Logistic\\Interpretable Odds Ratios};
\node[block, align=center] (n5) at (12.8,0) {5. Model B: CART Tree\\Nonlinear Splitting Rules};
\node[block, align=center] (n6) at (12.8,-2.2) {6. Model C: k-NN\\Instance-Based Memory};
\node[block, align=center] (n7) at (9.6,-2.2) {7. Stratified 10-Fold CV\\Fold Leakage Prevention};
\node[block, align=center] (n8) at (6.4,-2.2) {8. Metric Evaluation\\Sensitivity, F1, ROC-AUC};
\node[block, align=center] (n9) at (3.2,-2.2) {9. Brier Calibration\\Reliability Curve Fitting};
\node[block, align=center] (n10) at (0,-2.2) {10. Clinical Decision Tool\\Early Risk Stratification};
\draw[arr] (n1) -- (n2);
\draw[arr] (n2) -- (n3);
\draw[arr] (n3) -- (n4);
\draw[arr] (n4) -- (n5);
\draw[arr] (n5) -- (n6);
\draw[arr] (n6) -- (n7);
\draw[arr] (n7) -- (n8);
\draw[arr] (n8) -- (n9);
\draw[arr] (n9) -- (n10);
\end{tikzpicture}}
\caption{System Architecture and Algorithmic Workflow for Comparative Study of Classification Algorithms for Early Disease Risk Prediction.}
\label{fig:pipeline}
\end{figure}

\subsection{Component Descriptions}
\begin{enumerate}
  \item \textbf{Data Ingestion and Cleansing:} Ingests raw multidimensional observational tensors, executing median imputation for missing physiological or telemetry fields and stripping low-variance noise attributes.
  \item \textbf{Transformation and Scaling:} Applies robust standardization ($z = (x - \mu) / \sigma$) exclusively across training folds to prevent target and distribution leakage into evaluation test splits.
  \item \textbf{Core Algorithmic Estimation:} Executes parameter optimization utilizing analytical normal equations, iteratively reweighted least squares, or gradient-based solvers.
  \item \textbf{Hyperparameter Optimization:} Evaluates cross-validation performance across logarithmic candidate grids, applying the one-standard-error rule to select the most parsimonious model.
  \item \textbf{Validation and Audit Engine:} Generates out-of-fold confusion matrices, calibration plots, and variance stability profiles across multi-seed Monte Carlo splits.
\end{enumerate}
